\section{はじめに} \label{sec1}
% Llama3.1~\cite{Meta24}やQwen3~\cite{qwen25}のような大規模言語モデル（Large Language Model, LLM）は，多様なタスクで高い性能を発揮している~\cite{zhao25}．
% Llama3.1~\cite{Meta24}やQwen3~\cite{qwen25}のような
Large Language Model（LLM）は多様なタスクで高い性能を発揮している~\cite{Zhao25}が，LLMの出力は学習データ内の知識に依存するため，未学習の知識に対して誤情報を生成する可能性がある~\cite{Ji23}．
また，日々更新される情報のすべてを網羅して学習することは困難である~\cite{Wang25}．
これらの課題に対する手法として，外部知識を検索し，その結果をLLMのコンテキストとして提示するRetrieval-Augmented Generation（RAG）~\cite{Lewis20}が提案されている．
RAGにより，LLMは最新かつ正確な知識を反映した応答が可能となる．

RAGの活用例の一つに，質問応答（Question Answering, QA）タスクがある．
RAGを用いたQAに関する初期の研究では，単一文書内に解答が存在するシングルホップ型のQAタスク~\cite{rajpurkar16}に焦点が当てられてきた~\cite{Shi24}．
しかし，複数文書間の関連付けや段階的推論を要する，マルチホップ型のQAタスクも多数存在する~\cite{yang18,ho20,trivedi22}．
例えば，「アインシュタインが相対性理論を発表したときのアメリカ大統領は誰ですか？」という質問は，発表時期と在任大統領を照合する二段階推論を必要とする．
このような複雑な質問に対応するため，LightRAG~\cite{Guo25}やGraphRAG~\cite{edge25}のようなマルチホップ型のQAタスクに焦点が当てられたRAGが提案されている．

一方で，これらの手法には，複数回の検索や段階的推論を必要とするため，推論時間が増大するという課題がある．
実運用において，入力される質問の全てが複雑であるとは限らない．
頻出する質問の中には単純かつ直接的なものも多く，例えば「アインシュタインの誕生日はいつですか？」のような質問は，外部検索を用いずともLLMの内部知識のみで回答できる場合がある~\cite{Mallen23}．
このような単純な質問に対しては，不要な検索や推論はコストを増大させるだけである．
%すなわち，複雑な質問には高精度なRAGが有効である一方，単純な質問に対しては不要な検索や推論がコストを増大させる．
しかし，複雑な質問を単純なRAG構成で処理した場合，回答精度が低下する可能性がある~\cite{LingjiaoChen24}．

% 研究課題
以上の背景から，質問の特徴に応じてRAGを動的に調整する「適応的なRAG」に関する研究が進展している~\cite{Asai23,fu24,Tang25,Rishi25,zhang2025adagcrag}．
% 具体的には，RAGのプロセス途中に検索の有無や関連文書の選択をLLMが調整する手法 ~\cite{Asai23}，オンライン強化学習によりハイパーパラメータや検索方法を最適化する手法~\cite{fu24,Tang25} ，質問に応じてファインチューニング済みモデルが検索方法を調整する手法~\cite{Rishi25}やハイパーパラメータを調整する手法~\cite{zhang2025adagcrag}も提案されている．
しかし，既存の適応的なRAGには，主に三つの課題がある．
    第一に，LLMの自己判定やプランニングを利用するエージェント指向の調整手法（Agentic RAG等）~\cite{Asai23}は，RAG構成の
決定にLLM自身の多段推論やオンライン探索を必要とする．そのため，回答精度が高まりやすい一方で，判断ごとに大きなLLM推論オー
バーヘッド（実行遅延やAPIコスト）が発生し，かつチャンクサイズや重複度といった複数のミクロなハイパーパラメータを総合的か
つ低コストに調整することが困難である．
    第二に，オンライン強化学習を用いた手法~\cite{fu24,Tang25}は，稼働中の環境での継続的な試行錯誤と学習を伴うため，実運用時のインフラコストや学習の不確実性が高まりやすい．
    第三に，質問や文書の特徴に基づく適応的アプローチ~\cite{Rishi25,zhang2025adagcrag}は存在するが，それらの多くはマクロな処理パラダイムの切り替え（例: 検索有無やモデルの動的ルーティング）や特定のグラフベースRAGパラメータの制御に特化しており，多様なRAG方式（Graph / Non-Graph）を横断し，かつチャンクサイズや重複度といったミクロなハイパーパラメータまでを含めた大規模な離散構成空間（全96通り）を一括して順位付けするアプローチは十分に議論されていない．また，対立する多目的指標（適合率，再現率，F1値，実行時間）を実行時にユーザ定義の優先度（重み $W$）に応じて動的に制御できる統合的なフレームワークも限られている．

% 研究目的
これらの課題に対し，本研究はLLMやオンラインでの継続的な学習・探索に依存せず，事前に収集した実測データから多様なRAG構成を学習するオフライン学習を採用する．
推論時に追加のLLM推論を必要とする調整ではなく，あらかじめ定義したRAG構成空間の候補に対して，精度指標とコスト指標それぞれのランキングを学習し，推論時は学習済みランキングモデル（Ranking Model, RM）の推論と任意の重みでRAG構成を選択する．
これにより，オンライン探索を伴うエージェント指向の調整手法と比較して，推論コストと応答遅延（ミリ秒オーダーの決定）を大幅に抑えつつ，複数のRAGとハイパーパラメータを含む広い構成空間を一括して扱うことを目指す．
また，オンライン学習を前提としないため，第二の課題である運用時の継続的な学習コストも回避する．
さらに，質問側の特徴に加えて文書側の表層的特徴および意味的特徴を用い，精度指標とコスト指標を統合した多目的な観点から学習することで，第三の課題に対応する．

% 手法概説
本手法では，まず質問と文書から統計量に基づく表層的特徴および埋め込み表現に基づく意味的特徴を抽出する．
学習時には，多様な RAG 構成に対する実測性能を教師データとしてRMを学習する．
% 実測性能（適合率，再現率，$F_1$ 値，時間）
推論時には，このRMが各質問に対して候補構成を順位付けし，上位の構成を選択することで，低コストや高精度のRAG構成を選択する．
さらに，精度重視，バランス重視，コスト重視といった要求に応じて重み付けを切り替え，複数のランキング結果を使い分けることで，ユーザの多様な要求に対応する．

% 実験概説
評価実験では，QAタスクを対象として，RMの性能と，本手法により選択されたRAG構成を用いた最終的な応答を評価する．
%実験の結果，ランキングモデルはランダムと比較して高い精度で順位を予測できており，特に時間の予測において優れた性能を示した．
%また，本手法によって択されたRAG構成は，質問に対して同一の RAG 構成を適用した場合と比較して， $F_1$値 を改善しつつ時間を短縮した．
%加えて，精度とコストのトレードオフを運用方針に応じて調整できることを確認した．
%さらに，ユーザ選択を含む適応的なRAG構成選択を用いた実験では，本手法の実用性が示され，特に時間重視の設定ではオラクルと同等の性能に到達した．

本研究の貢献は以下の通りである．
\begin{itemize}
\item RAG構成情報（生成モデルや検索方式等）が順位付けの基盤となる一方，質問・文書の表層的および意味的特徴量を包括的に統合することで，構成単独（NDCG@1=66.84\%）を上回る最高精度（NDCG@1=68.52\%，Hit@5=65.00\%）のランキングが可能となることを実験的に明らかにした．
\item 精度（適合率，再現率，$F_1$値）とコスト（実行時間）を統合的に扱う多目的ランキング手法を導入し，ユーザの優先順位や入力状況に応じたRAG構成の適応的な自動選択を実現した．
\end{itemize}

% \item 質問および文書の表層的特徴と意味的特徴を統合したRAG構成選択用の特徴量を設計し，それらを考慮したランキングモデルを構築した．
% \item 精度（適合率，再現率，$F_1$値）とコスト（時間）を統合的に評価する多目的ランキング手法を導入し，ユーザの要求に応じて適応的なRAG構成を選択する手法を提案した．
本論文の構成は次の通りである．
第~\ref{sec2}では本研究の関連研究を概説し，第~\ref{sec3}では提案手法について詳述する.
第~\ref{sec4}では評価実験の方法について述べ，第~\ref{sec5}では評価実験結果および考察を述べる．
最後に，第~\ref{sec6}では本研究の結論と今後の課題を提示する．

なお，本論文は，穴口らの研究成果~\cite{anaguchi26}に対し，以下の具体的な追加・拡張を行ったものである．
第一に，近接研究（CARROT，uRAG，RAGRouter，AdaGCRAG）との関係性を整理し，手法のポジショニングを明確化した．
第二に，既存の代表的な適応的RAGベースライン手法であるAdaptive-RAGを評価フレームワークに組み込み，精度とコストのParetoフロンティアの観点から性能評価を行った．
第三に，自動評価指標（BERTScore）に加えて，標準的QA評価指標であるExact Match (EM)およびToken-level $F_1$値を用いて提案手法を多角的に再評価した．
第四に，提案手法とベースライン間の差について，Wilcoxon符号付き順位検定による有意差検定を行い，回答精度の統計的優位性を立証した．
最後に，上位K件提示モデルにおける構成順次試行コストを「累積実行時間」として計上し，より実用的なコスト効率性を実証した．